\documentclass[12pt]{article}
\usepackage[UTF8]{inputenc}
\usepackage{thaisupp}
\usepackage{geometry}
\geometry{a4paper, margin=2.5cm}
\begin{document}
\begin{center}
{\LARGE\textbf{ฟิสิกส์ ม.ปลาย — Part 1}}\\[0.3em]
{\Large\textbf{Lesson 10: เสียง}}\\[0.5em]
\rule{0.8\linewidth}{0.5pt}
\end{center}
\section*{คำถาม In-Class}
\begin{enumerate}
\item \textbf{เสียงเดินทางในอากาศที่อุณหภูมิ 25°C ด้วยอัตราเร็วประมาณเท่าใด}
\begin{enumerate}
\item ก) 300 m/s
\item ข) 340 m/s
\item ค) 380 m/s
\item ง) 400 m/s
\end{enumerate}
\item \textbf{เสียงที่มนุษย์ได้ยินมีความถี่ในช่วงใด}
\begin{enumerate}
\item ก) 20 Hz - 20,000 Hz
\item ข) 2 Hz - 2,000 Hz
\item ค) 200 Hz - 20,000 Hz
\item ง) 20 kHz - 20,000 kHz
\end{enumerate}
\item \textbf{เมื่อรถพยุมยางเสียงดังเข้าหาผู้สังเกตที่อยู่กับที่ ผู้สังเกตจะได้ยินเสียงอย่างไร}
\begin{enumerate}
\item ก) เสียงทุ้มลง
\item ข) เสียงแหลมสูงขึ้น
\item ค) เสียงเหมือนเดิม
\item ง) เสียงดังขึ้นแต่ระดับเสียงเหมือนเดิม
\end{enumerate}
\item \textbf{ระดับเสียง 60 dB มีความเข้มเท่ากับกี่เท่าของระดับเสียง 0 dB}
\begin{enumerate}
\item ก) 60 เท่า
\item ข) 600 เท่า
\item ค) 10⁶ เท่า
\item ง) 10⁸ เท่า
\end{enumerate}
\item \textbf{แหล่งเสียงความถี่ 500 Hz เคลื่อนที่ด้วยความเร็ว 30 m/s เข้าหาผู้สังเกตที่อยู่กับที่ ความถี่ที่ผู้สังเกตได้ยินเป็นเท่าใด (v = 340 m/s)}
\begin{enumerate}
\item ก) 444 Hz
\item ข) 500 Hz
\item ค) 556 Hz
\item ง) 600 Hz
\end{enumerate}
\item \textbf{คลื่นเสียงความถี่ 256 Hz มีความยาวคลื่นเท่าใดในอากาศ (v = 340 m/s)}
\begin{enumerate}
\item ก) 0.75 m
\item ข) 1.33 m
\item ค) 2.66 m
\item ง) 3.75 m
\end{enumerate}
\item \textbf{ข้อใดต่อไปนี้ไม่ใช่ปัจจัยที่มีผลต่ออัตราเร็วเสียงในอากาศ}
\begin{enumerate}
\item ก) อุณหภูมิ
\item ข) ความดัน
\item ค) ความชื้น
\item ง) แหล่งกำเนิดเสียง
\end{enumerate}
\item \textbf{เครื่องดนตรีชนิดใดที่ให้เสียงความถี่พื้นฐานสูงสุด}
\begin{enumerate}
\item ก) กีตาร์
\item ข) ไวโอลิน
\item ค) ไพโอท
\item ง) เปียโน
\end{enumerate}
\item \textbf{รถไฟเคลื่อนที่ด้วยความเร็ว 40 m/s ส่งเสียงนกหวีดความถี่ 400 Hz อัตราเร็วเสียง 340 m/s ความถี่ที่คนบนรถไฟได้ยินจะเป็นอย่างไร}
\begin{enumerate}
\item ก) สูงขึ้น
\item ข) ต่ำลง
\item ค) เท่าเดิม
\item ง) ขึ้นอยู่กับทิศทางลม
\end{enumerate}
\item \textbf{ระดับเสียงเพิ่มขึ้นจาก 40 dB เป็น 60 dB ความเข้มเสียงเพิ่มขึ้นกี่เท่า}
\begin{enumerate}
\item ก) 2 เท่า
\item ข) 20 เท่า
\item ค) 100 เท่า
\item ง) 1000 เท่า
\end{enumerate}
\item \textbf{เสียงเดินทางในของเหลวเร็วกว่าในแก๊ส เพราะเหตุใด}
\begin{enumerate}
\item ก) ของเหลวมีความหนาแน่นมากกว่า
\item ข) โมเลกุลในของเหลวอยู่ชิดกันมากกว่า ส่งผ่านการสั่นได้เร็วกว่า
\item ค) ของเหลวมีอุณหภูมิสูงกว่า
\item ง) ของเหลวมีความดันต่ำกว่า
\end{enumerate}
\item \textbf{รถพยุมยางเสียงดังความถี่ 300 Hz วิ่งเข้าหาผู้สังเกตด้วยความเร็ว 25 m/s ห่างจากผู้สังเกต 500 m ความถี่ที่ผู้สังเกตได้ยินขณะรถเข้าใกล้เป็นเท่าใด (v = 340 m/s)}
\begin{enumerate}
\item ก) 322 Hz
\item ข) 325 Hz
\item ค) 525 Hz
\item ง) 600 Hz
\end{enumerate}
\item \textbf{หน่วยเดซิเบล (dB) ใช้วัดอะไร}
\begin{enumerate}
\item ก) ความถี่เสียง
\item ข) ความยาวคลื่นเสียง
\item ค) ระดับเสียง (ความดัง)
\item ง) อัตราเร็วเสียง
\end{enumerate}
\item \textbf{เสียงอัลตราซาวด์มีความถี่สูงกว่าเท่าใด}
\begin{enumerate}
\item ก) 200 Hz
\item ข) 2,000 Hz
\item ค) 20,000 Hz
\item ง) 200,000 Hz
\end{enumerate}
\item \textbf{เมื่อแหล่งเสียงเคลื่อนที่ห่างออกจากผู้สังเกตที่อยู่กับที่ ความถี่ที่ได้ยินจะเป็นอย่างไร}
\begin{enumerate}
\item ก) สูงขึ้น
\item ข) ต่ำลง
\item ค) เท่าเดิม
\item ง) เป็นศูนย์
\end{enumerate}
\item \textbf{เสียงที่มนุษย์ได้ยินเกิดจากการสั่นของวัตถุในอากาศด้วยความถี่ประมาณเท่าใดจึงจะได้ยิน}
\begin{enumerate}
\item ก) 0.02 Hz - 20 Hz
\item ข) 20 Hz - 20,000 Hz
\item ค) 20,000 Hz - 200,000 Hz
\item ง) 200 Hz - 2,000 Hz
\end{enumerate}
\item \textbf{ระดับเสียง 80 dB มีความเข้มเป็นกี่เท่าของ 50 dB}
\begin{enumerate}
\item ก) 30 เท่า
\item ข) 100 เท่า
\item ค) 1000 เท่า
\item ง) 10,000 เท่า
\end{enumerate}
\item \textbf{เสียงในข้อใดเดินทางได้เร็วที่สุด}
\begin{enumerate}
\item ก) ในอากาศ
\item ข) ในน้ำ
\item ค) ในเหล็ก
\item ง) ในไม้
\end{enumerate}
\item \textbf{แหล่งเสียงความถี่ 600 Hz เคลื่อนที่ด้วยความเร็ว 20 m/s เข้าหาผู้สังเกตที่อยู่กับที่ ความถี่ที่ได้ยินเป็นเท่าใด (v = 340 m/s)}
\begin{enumerate}
\item ก) 565 Hz
\item ข) 636 Hz
\item ค) 680 Hz
\item ง) 720 Hz
\end{enumerate}
\item \textbf{เสียงดัง 120 dB สามารถทำลายการได้ยินได้หรือไม่}
\begin{enumerate}
\item ก) ไม่ ต้องดังกว่า 200 dB
\item ข) ได้ เพราะเป็นระดับเสียงที่เจ็บปวด
\item ค) ได้เฉพาะเด็ก
\item ง) ไม่เคยทำลายการได้ยิน
\end{enumerate}
\item \textbf{ความถี่ของเสียงกำทินหนึ่ง 512 Hz เดินทางในอากาศด้วยอัตราเร็ว 340 m/s ความยาวคลื่นเป็นเท่าใด}
\begin{enumerate}
\item ก) 0.5 m
\item ข) 0.67 m
\item ค) 1.5 m
\item ง) 2 m
\end{enumerate}
\item \textbf{รถพยุมยางเสียงดัง 500 Hz วิ่งด้วยความเร็ว 30 m/s ผ่านผู้สังเกตที่ยืนอยู่ริมถนน ความถี่ที่ผู้สังเกตได้ยินเมื่อรถเข้าใกล้และเมื่อรถผ่านพ้นต่างกันเท่าใด (v = 340 m/s)}
\begin{enumerate}
\item ก) ประมาณ 90 Hz
\item ข) ประมาณ 46 Hz
\item ค) ประมาณ 22 Hz
\item ง) ประมาณ 10 Hz
\end{enumerate}
\item \textbf{อัตราเร็วเสียงในอากาศขึ้นอยู่กับปัจจัยใดเป็นหลัก}
\begin{enumerate}
\item ก) อุณหภูมิ
\item ข) ความดันบรรยากาศ
\item ค) ความชื้นสัมพัทธ์
\item ง) มวลอากาศ
\end{enumerate}
\item \textbf{เครื่องขยายเสียงเพิ่มกำลังเสียงจาก 1 W เป็น 100 W ระดับเสียงเพิ่มขึ้นกี่ dB}
\begin{enumerate}
\item ก) 10 dB
\item ข) 20 dB
\item ค) 100 dB
\item ง) 40 dB
\end{enumerate}
\item \textbf{ผู้สังเกตเคลื่อนที่ด้วยความเร็ว 10 m/s เข้าหาแหล่งเสียงที่อยู่กับที่ความถี่ 300 Hz ความถี่ที่ได้ยินเป็นเท่าใด (v = 340 m/s)}
\begin{enumerate}
\item ก) 280 Hz
\item ข) 300 Hz
\item ค) 309 Hz
\item ง) 320 Hz
\end{enumerate}
\end{enumerate}
\end{document}